\subsection{Tecnologías y enfoques}

Con el objetivo de abordar las limitaciones de los sistemas 
centralizados, han surgido diversas tecnologías descentralizadas 
orientadas a mejorar la disponibilidad, integridad y verificabilidad 
de los datos. Entre ellas, destacan los sistemas de almacenamiento 
distribuido basados en contenido y las redes blockchain, 
que representan dos enfoques complementarios para la gestión 
de información en entornos sin confianza.

El sistema \textit{InterPlanetary File System} (IPFS) propone un 
modelo de almacenamiento distribuido en el que los ficheros se 
identifican mediante hashes criptográficos derivados de su contenido, 
en lugar de mediante ubicaciones físicas~\cite{benet2014ipfs}.
Este enfoque permite garantizar la integridad de los datos y 
facilita su distribución eficiente entre nodos de la red. 
Además, IPFS mejora la disponibilidad de los ficheros, 
ya que estos pueden ser replicados y servidos desde múltiples nodos, 
eliminando la dependencia de un servidor central.

Por otro lado, las tecnologías blockchain introducen un modelo de 
registro distribuido e inmutable que permite almacenar información 
de forma verificable y resistente a manipulaciones. En particular,
plataformas como Ethereum permiten ejecutar contratos inteligentes
y registrar datos de manera transparente en una red 
descentralizada~\cite{wood2014ethereum}. 
Sin embargo, debido a las limitaciones de escalabilidad 
y al coste asociado al almacenamiento en blockchain,
no resulta eficiente almacenar directamente grandes 
volúmenes de datos en estas redes.

Como consecuencia, han surgido enfoques híbridos que combinan
ambas tecnologías con el fin de aprovechar sus ventajas respectivas.
En estos sistemas, los ficheros se almacenan en redes
distribuidas como IPFS, mientras que sus identificadores (hashes)
se registran en blockchain para garantizar su autenticidad 
e inmutabilidad. Este modelo permite desacoplar el 
almacenamiento de datos del mecanismo de verificación,
logrando un equilibrio entre eficiencia, coste y seguridad.

Este trabajo adopta este enfoque híbrido mediante el 
desarrollo de un sistema que integra almacenamiento descentralizado 
basado en IPFS con un mecanismo de verificación implementado 
sobre una red blockchain local. Este planteamiento permite evaluar 
de forma controlada el comportamiento de ambas tecnologías y 
analizar sus compromisos en términos de rendimiento y funcionalidad.